\documentclass[10pt,a4paper]{amsart}
\usepackage[margin=19mm]{geometry}
\usepackage{amsmath,amssymb,mathtools}
\usepackage[scr=rsfs,cal=euler]{mathalfa}
\usepackage{xfrac}
\usepackage[unicode,hidelinks,bookmarks=false]{hyperref}
\newcommand{\ie}{i.e.\ }
\newcommand{\cf}{cf.\ }
\newcommand{\resp}{respectively\ }
\newcommand{\Subsection}{\subsection}
\providecommand{\nobreakdash}{\nobreak\hbox{-}}
\setlength{\emergencystretch}{2em}
\multlinegap0pt
\usepackage{siunitx}
\usepackage{siunitx}
\usepackage{siunitx}
\usepackage{siunitx}
\usepackage{mathtools}
\usepackage{tcolorbox}
\usepackage{tikz}
\usepackage{csquotes}
\usepackage{siunitx}
\usepackage{float}
\usepackage{amssymb}
\usepackage{algorithm}\usepackage{algpseudocode}
\usepackage[shortlabels]{enumitem}
\usepackage{xcolor}
\newtheorem{lemma}{Lemma}
\newcommand{\R}{\bb R}
\newcommand{\bb}{\mathbb}
\newcommand{\eqdef}{:=}
\newcommand{\op}{\operatorname{op}}
\newcommand{\opnorm}[1]{\left\| #1 \right\|_{\op}}      % norm 
\providecommand{\paragraph}{}
\newcommand{\rb}[1]{\left(#1\right)}
\newcommand{\sqn}[1]{\left\| #1 \right\|^2}
\title{SILAGE: Memory-Efficient, Full-Gradient-Free Nonconvex Optimization for Nested Finite Sums}
\author{Igor Sokolov, Laurent Condat, Peter Richt\'arik}
\date{Source: arXiv:2606.15832v2 (CC BY 4.0)}
\begin{document}
\maketitle

\noindent\emph{Source and licence.} SILAGE: Memory-Efficient, Full-Gradient-Free Nonconvex Optimization for Nested Finite Sums. Version arXiv:2606.15832v2 (CC BY 4.0), distributed by arXiv under the Creative Commons Attribution 4.0 International licence, http://creativecommons.org/licenses/by/4.0/, as recorded in the arXiv licence metadata for this exact version and shown on the abstract page https://arxiv.org/abs/2606.15832v2. Full licence notice: \texttt{licenses/arxiv-2606.15832-CC-BY-4.0.txt}.\par\smallskip

\noindent\textit{Editorial scope.} Counted unit: the article's lemma lem:hessian\_characterization\_delta (source0.tex lines 5305--5380). The article's complete original proof of it is delivered with it (source0.tex lines 5382--5569, one \texttt{proof} environment of 188 lines). No mathematics is written, repaired or re-derived: the statement and the proof are the article's own lines, in the article's own notation, and no retained line is substituted or re-flowed.  Retained uncounted as the source-local context the proof needs: lines 5295--5300.  Nothing was omitted from any retained span, so the package carries no omission record.  No retained line contains a citation, so the wrapper carries no bibliography and no bibliography entry.  No retained line contains a figure environment, an image-inclusion command or drawing code, so no image is delivered and the package declares no asset dependency.  Editorial formatting only, in addition: an AMS \texttt{amsart} preamble that restates the article's own declarations and macros used by the retained lines, namely the environments \texttt{lemma} on separate counters, together with the standard packages those lines need.  The retained environments print in this document as Lemma 1 in order of appearance, whereas the article prints the same items with the numbering of its own full document.  The complete native source of the article as distributed in the arXiv e-print for this exact version is bundled unchanged as \texttt{sources/202894/source0.tex}.
\begin{align}
\label{eq:exp_delta1}
\frac{1}{n}\sum_{i=1}^n \sqn{\nabla f_i(x)-\nabla f(x)-\nabla f_i(y)+\nabla f(y)} &\le \delta_1^2\sqn{x-y}, \\
\label{eq:exp_delta2}
\frac{1}{nm}\sum_{i=1}^n\sum_{j=1}^m \sqn{\nabla f_{i,j}(x)-\nabla f_i(x)-\nabla f_{i,j}(y)+\nabla f_i(y)} &\le \delta_2^2\sqn{x-y}.
\end{align}

\begin{lemma}[Hessian characterization of $\delta_1$- and $\delta_2$-similarity]
\label{lem:hessian_characterization_delta}
Let $f(x)\eqdef \frac{1}{n}\sum_{i=1}^n f_i(x)$ and
$f_i(x)\eqdef \frac{1}{m}\sum_{j=1}^m f_{i,j}(x)$.

\smallskip
\noindent
\textbf{(i) $\delta_1$-similarity.}
Assume that $f_i$ is twice continuously differentiable for all $i\in[n]$.
For each $x\in\R^d$, define the linear operator
$\mathcal{A}_1(x):\R^d\to(\R^d)^n$ by
\begin{align}
    \mathcal{A}_1(x)s
    \eqdef
    \frac{1}{\sqrt n}
    \big(
    (\nabla^2 f_i(x)-\nabla^2 f(x))s
    \big)_{i=1}^n .
    \label{eq:A1_definition}
\end{align}
Then \eqref{eq:exp_delta1} holds if and only if
\begin{align}
    \opnorm{\mathcal{A}_1(x)}\le \delta_1,
    \qquad
    \forall x\in\R^d .
    \label{eq:hess_delta1_stacked_opnorm}
\end{align}
Equivalently,
\begin{align}
    \lambda_{\max}\rb{
    \frac1n\sum_{i=1}^n
    \rb{\nabla^2 f_i(x)-\nabla^2 f(x)}^\top
    \rb{\nabla^2 f_i(x)-\nabla^2 f(x)}
    }
    \le \delta_1^2,
    \qquad
    \forall x\in\R^d .
    \label{eq:hess_delta1_spectral}
\end{align}

\smallskip
\noindent
\textbf{(ii) $\delta_2$-similarity.}
Assume that $f_{i,j}$ is twice continuously differentiable for all
$i\in[n]$ and $j\in[m]$.
For each $x\in\R^d$, define the linear operator
$\mathcal{A}_2(x):\R^d\to(\R^d)^{nm}$ by
\begin{align}
    \mathcal{A}_2(x)s
    \eqdef
    \frac{1}{\sqrt{nm}}
    \big(
    (\nabla^2 f_{i,j}(x)-\nabla^2 f_i(x))s
    \big)_{i\in[n],\,j\in[m]} .
    \label{eq:A2_definition}
\end{align}
Then \eqref{eq:exp_delta2} holds if and only if
\begin{align}
    \opnorm{\mathcal{A}_2(x)}\le \delta_2,
    \qquad
    \forall x\in\R^d .
    \label{eq:hess_delta2_stacked_opnorm}
\end{align}
Equivalently,
\begin{align}
    \lambda_{\max}\rb{
    \frac1{nm}\sum_{i=1}^n\sum_{j=1}^m
    \rb{\nabla^2 f_{i,j}(x)-\nabla^2 f_i(x)}^\top
    \rb{\nabla^2 f_{i,j}(x)-\nabla^2 f_i(x)}
    }
    \le \delta_2^2,
    \qquad
    \forall x\in\R^d .
    \label{eq:hess_delta2_spectral}
\end{align}
\end{lemma}

\begin{proof}
We prove the two equivalences separately.

\paragraph{(i) $\delta_1$-similarity.}
Define the centered gradients
\begin{align}
    G_i(x)\eqdef \nabla f_i(x)-\nabla f(x).
    \notag
\end{align}
Then \eqref{eq:exp_delta1} is equivalent to
\begin{align}
    \frac1n\sum_{i=1}^n
    \sqn{G_i(x)-G_i(y)}
    \le
    \delta_1^2\sqn{x-y},
    \qquad
    \forall x,y\in\R^d .
    \label{eq:delta1_G_form}
\end{align}

\smallskip
\noindent
\emph{($\Rightarrow$)}
Fix $x\in\R^d$, $s\in\R^d$, and $\alpha>0$, and set
$y=x+\alpha s$. Applying \eqref{eq:delta1_G_form} gives
\begin{align}
    \frac1n\sum_{i=1}^n
    \sqn{
    \frac{G_i(x+\alpha s)-G_i(x)}{\alpha}}
    \le
    \delta_1^2\sqn{s}.
    \label{eq:delta1_scaled_difference}
\end{align}
Since each $f_i$ is twice continuously differentiable,
\begin{align}
    \frac{G_i(x+\alpha s)-G_i(x)}{\alpha}
    \xrightarrow[\alpha\downarrow 0]{}
    \rb{\nabla^2 f_i(x)-\nabla^2 f(x)}s .
    \notag
\end{align}
Taking $\alpha\downarrow 0$ in \eqref{eq:delta1_scaled_difference} yields
\begin{align}
    \frac1n\sum_{i=1}^n
    \sqn{
    \rb{\nabla^2 f_i(x)-\nabla^2 f(x)}s}
    \le
    \delta_1^2\sqn{s},
    \qquad
    \forall x,s\in\R^d .
    \label{eq:delta1_directional_hessian}
\end{align}
By the definition of $\mathcal{A}_1(x)$, this is exactly
$\sqn{\mathcal{A}_1(x)s}\le \delta_1^2\sqn{s}$ for all $s$,
which is equivalent to $\opnorm{\mathcal{A}_1(x)}\le \delta_1$.

\smallskip
\noindent
\emph{($\Leftarrow$)}
Assume \eqref{eq:hess_delta1_stacked_opnorm}. Fix $x,y\in\R^d$
and set $s=y-x$. By the fundamental theorem of calculus,
\begin{align}
    G_i(y)-G_i(x)
    =
    \int_0^1
    \rb{\nabla^2 f_i(x+\tau s)-\nabla^2 f(x+\tau s)}s
    \,d\tau .
    \notag
\end{align}
Stacking the vectors and using Jensen's inequality in the product space
$(\R^d)^n$, we obtain
\begin{align}
    \frac1n\sum_{i=1}^n \sqn{G_i(y)-G_i(x)}
    &=
    \sqn{
    \int_0^1
    \mathcal{A}_1(x+\tau s)s
    \,d\tau}
    \notag
    \\
    &\le
    \int_0^1
    \sqn{\mathcal{A}_1(x+\tau s)s}
    \,d\tau
    \notag
    \\
    &\le
    \int_0^1
    \delta_1^2\sqn{s}
    \,d\tau
    =
    \delta_1^2\sqn{y-x}.
    \notag
\end{align}
This is \eqref{eq:exp_delta1}.

\paragraph{(ii) $\delta_2$-similarity.}
Define the within-group centered gradients
\begin{align}
    G_{i,j}(x)\eqdef \nabla f_{i,j}(x)-\nabla f_i(x).
    \notag
\end{align}
Then \eqref{eq:exp_delta2} is equivalent to
\begin{align}
    \frac1{nm}\sum_{i=1}^n\sum_{j=1}^m
    \sqn{G_{i,j}(x)-G_{i,j}(y)}
    \le
    \delta_2^2\sqn{x-y},
    \qquad
    \forall x,y\in\R^d .
    \label{eq:delta2_H_form}
\end{align}

\smallskip
\noindent
\emph{($\Rightarrow$)}
Fix $x\in\R^d$, $s\in\R^d$, and $\alpha>0$, and set
$y=x+\alpha s$. Applying \eqref{eq:delta2_H_form} gives
\begin{align}
    \frac1{nm}\sum_{i=1}^n\sum_{j=1}^m
    \sqn{
    \frac{G_{i,j}(x+\alpha s)-G_{i,j}(x)}{\alpha}}
    \le
    \delta_2^2\sqn{s}.
    \label{eq:delta2_scaled_difference}
\end{align}
Since the functions $f_{i,j}$ are twice continuously differentiable,
\begin{align}
    \frac{G_{i,j}(x+\alpha s)-G_{i,j}(x)}{\alpha}
    \xrightarrow[\alpha\downarrow 0]{}
    \rb{\nabla^2 f_{i,j}(x)-\nabla^2 f_i(x)}s .
    \notag
\end{align}
Taking $\alpha\downarrow 0$ in \eqref{eq:delta2_scaled_difference} yields
\begin{align}
    \frac1{nm}\sum_{i=1}^n\sum_{j=1}^m
    \sqn{
    \rb{\nabla^2 f_{i,j}(x)-\nabla^2 f_i(x)}s}
    \le
    \delta_2^2\sqn{s},
    \qquad
    \forall x,s\in\R^d .
    \label{eq:delta2_directional_hessian}
\end{align}
By the definition of $\mathcal{A}_2(x)$, this is exactly
$\sqn{\mathcal{A}_2(x)s}\le \delta_2^2\sqn{s}$ for all $s$,
which is equivalent to $\opnorm{\mathcal{A}_2(x)}\le \delta_2$.

\smallskip
\noindent
\emph{($\Leftarrow$)}
Assume \eqref{eq:hess_delta2_stacked_opnorm}. Fix $x,y\in\R^d$
and set $s=y-x$. By the fundamental theorem of calculus,
\begin{align}
    G_{i,j}(y)-G_{i,j}(x)
    =
    \int_0^1
    \rb{\nabla^2 f_{i,j}(x+\tau s)-\nabla^2 f_i(x+\tau s)}s
    \,d\tau .
    \notag
\end{align}
Stacking the vectors and using Jensen's inequality in the product space
$(\R^d)^{nm}$, we obtain
\begin{align}
    \frac1{nm}\sum_{i=1}^n\sum_{j=1}^m
    \sqn{G_{i,j}(y)-G_{i,j}(x)}
    &=
    \sqn{
    \int_0^1
    \mathcal{A}_2(x+\tau s)s
    \,d\tau}
    \notag
    \\
    &\le
    \int_0^1
    \sqn{\mathcal{A}_2(x+\tau s)s}
    \,d\tau
    \notag
    \\
    &\le
    \int_0^1
    \delta_2^2\sqn{s}
    \,d\tau
    =
    \delta_2^2\sqn{y-x}.
    \notag
\end{align}
This is \eqref{eq:exp_delta2}.
\end{proof}

\end{document}
